% !TeX program = pdflatex
% !TeX spellcheck = es_ES
%=======================================================================
%  INFORME DE LABORATORIO — FORMATO IEEE (clase IEEEtran)
%  Tema: Simulacion de un automatismo electrico (diagrama ladder) en CADe SIMU
%
%  COMO COMPILAR EN OVERLEAF:
%    Menu > Settings > Compiler  ->  pdfLaTeX
%    Menu > Settings > TeX Live version -> la mas reciente
%  El logo va en figuras/logo-universidad.png (ver README.md)
%=======================================================================
\documentclass[conference,a4paper,10pt]{IEEEtran}

%----------------------- Idioma y codificacion -------------------------
\usepackage[utf8]{inputenc}
\usepackage[T1]{fontenc}
\usepackage[spanish,es-tabla,es-nodecimaldot]{babel}
\usepackage{mathptmx}              % tipografia Times, la del estilo IEEE

%----------------------------- Utilidades ------------------------------
\usepackage{graphicx}
\usepackage{amsmath,amssymb}
\usepackage{booktabs}
\usepackage{multirow}
\usepackage{array}
\usepackage{enumitem}
\usepackage{tikz}                  % para el diagrama de tiempos
\usepackage{rotating}              % permite figuras giradas a pagina completa
\usepackage{cite}                  % agrupa citas: [1]-[3]
\usepackage{eso-pic}               % coloca el logo en el margen superior
\usepackage[hidelinks,breaklinks=true]{hyperref}

%=======================================================================
%  1. DATOS INSTITUCIONALES  ->  EDITA SOLO ESTE BLOQUE
%=======================================================================
\newcommand{\universidad}{UNIVERSIDAD XXXXXXXX}
\newcommand{\facultad}{Facultad de Ingeniería}
\newcommand{\programa}{Programa de Ingeniería Electrónica}
\newcommand{\asignatura}{Automatismos Eléctricos --- Grupo XX}
\newcommand{\fechainforme}{Ciudad, \today}

% Archivo del logo (CON extension). Si no existe, se dibuja un recuadro
% de reserva para que el documento compile igual.
\newcommand{\archivologo}{figuras/logo-universidad.png}

% Geometria del encabezado. NO hace falta cuadrar nada a mano: el bloque
% se mide al compilar y el titulo baja solo lo necesario para dejar
% \holguratitulo de aire, aunque cambies el alto del logo.
\newlength{\altologo}      \setlength{\altologo}{13mm}      % alto del logo
\newlength{\margensuplogo} \setlength{\margensuplogo}{6mm}  % borde del papel -> encabezado
\newlength{\margenlatlogo} \setlength{\margenlatlogo}{18mm} % margen lateral del encabezado
\newlength{\holguratitulo} \setlength{\holguratitulo}{5mm}  % aire entre encabezado y titulo

%=======================================================================
%  2. MAQUINARIA DEL ENCABEZADO  (no necesitas tocar nada de aqui)
%
%  Que hace: arma el bloque "logo + datos de la universidad", MIDE su
%  alto real, lo dibuja en el margen superior de la primera pagina y
%  calcula cuanto debe bajar el titulo para no solaparse con el.
%=======================================================================
\newsavebox{\cjlogo}
\newsavebox{\cjtexto}
\newsavebox{\cjencabezado}
\newlength{\anchoencabezado}
\newlength{\altoencabezado}
\newlength{\topetexto}         % borde superior del papel -> area de texto
\newlength{\bajadaencabezado}
\newlength{\altoprimeralinea}  % del borde del area de texto a la 1a linea base
\newlength{\auxenc}
\newlength{\alzadotitulo}      % alto aproximado de la primera linea del titulo
\setlength{\alzadotitulo}{6mm}
\setlength{\anchoencabezado}{\dimexpr\paperwidth-2\margenlatlogo\relax}

% Logo real, o recuadro de reserva si el archivo aun no existe
\newlength{\altocajareserva}
\setlength{\altocajareserva}{\altologo}
\addtolength{\altocajareserva}{-4pt}
\newcommand{\cajalogo}{%
  \expandafter\IfFileExists\expandafter{\archivologo}%
    {\includegraphics[height=\altologo]{\archivologo}}%
    {\setlength{\fboxsep}{2pt}%
     \fbox{\parbox[b][\altocajareserva][c]{32mm}%
       {\centering\scriptsize\textsf{LOGO DE LA\\UNIVERSIDAD}}}}%
}

% Texto institucional. Quita o agrega lineas con libertad: el alto se
% recalcula solo. Si te queda largo, cambia \scriptsize por \tiny.
\newcommand{\textoinstitucional}{%
  \raggedleft\scriptsize
  \textsc{\universidad}\\[1pt]
  \facultad\\[1pt]
  \programa\\[1pt]
  \asignatura\\[1pt]
  \fechainforme\par}

\newcommand{\encabezadoinstitucional}{%
  % -- distancia real del borde del papel al area de texto (la fija IEEEtran)
  \setlength{\topetexto}{1in}%
  \addtolength{\topetexto}{\voffset}%
  \addtolength{\topetexto}{\topmargin}%
  \addtolength{\topetexto}{\headheight}%
  \addtolength{\topetexto}{\headsep}%
  % -- 1) se mide el logo
  \sbox{\cjlogo}{\cajalogo}%
  % -- 2) el texto ocupa exactamente el ancho que sobra
  \setlength{\auxenc}{\anchoencabezado}%
  \addtolength{\auxenc}{-\wd\cjlogo}%
  \addtolength{\auxenc}{-8mm}%
  \ifdim\auxenc<40mm
    \typeout{^^JAVISO: el logo es demasiado ancho; reduce \string\altologo.^^J}%
    \setlength{\auxenc}{40mm}%
  \fi
  \sbox{\cjtexto}{\parbox{\auxenc}{\textoinstitucional}}%
  % -- 3) el alto del encabezado es el mayor de los dos
  \setlength{\altoencabezado}{\dimexpr\ht\cjlogo+\dp\cjlogo\relax}%
  \setlength{\auxenc}{\dimexpr\ht\cjtexto+\dp\cjtexto\relax}%
  \ifdim\auxenc>\altoencabezado \setlength{\altoencabezado}{\auxenc}\fi
  % -- 4) se componen logo y texto centrados verticalmente entre si
  \sbox{\cjencabezado}{%
    \makebox[\anchoencabezado]{%
      \parbox[b][\altoencabezado][c]{\wd\cjlogo}{\usebox{\cjlogo}}%
      \hfill
      \parbox[b][\altoencabezado][c]{\wd\cjtexto}{\usebox{\cjtexto}}}}%
  % -- 5) se dibuja en el margen superior, solo en esta pagina
  \setlength{\bajadaencabezado}{\margensuplogo}%
  \addtolength{\bajadaencabezado}{\ht\cjencabezado}%
  \AddToShipoutPictureBG*{%
    \AtPageUpperLeft{%
      \raisebox{-\bajadaencabezado}{%
        \makebox[\paperwidth][c]{\usebox{\cjencabezado}}}}}%
  % -- 6) cuanto hay que bajar la primera linea del titulo:
  %       (lo que el encabezado invade el area de texto) + holgura + alzado
  \setlength{\altoprimeralinea}{\margensuplogo}%
  \addtolength{\altoprimeralinea}{\altoencabezado}%
  \addtolength{\altoprimeralinea}{-\topetexto}%
  \addtolength{\altoprimeralinea}{\holguratitulo}%
  \addtolength{\altoprimeralinea}{\alzadotitulo}%
  \ifdim\altoprimeralinea<\z@ \setlength{\altoprimeralinea}{\z@}\fi
}

% Puntal invisible al inicio del titulo: baja su primera linea sin
% depender del modo en que IEEEtran componga el bloque de titulo.
\newcommand{\bajartitulo}{\noindent\rule{0pt}{\altoprimeralinea}}

% Para pasar el logo a la DERECHA y el texto a la IZQUIERDA, intercambia
% los dos \parbox del paso 4 (y cambia \raggedleft por \raggedright en
% \textoinstitucional).

% ---- Insercion de figuras tolerante a archivos faltantes --------------
% Uso: \imagen{<ancho>}{<ruta/archivo.ext>}
%   <ancho> NO debe pasar de \columnwidth en una figura de una columna
%   ni de \textwidth en una figure*: mas ancho que eso no cabe en el
%   papel, se sale por la derecha y \centering no lo puede centrar.
% La imagen se encaja dentro de <ancho> x \altomaximagen conservando la
% proporcion, asi que tampoco puede quedar mas alta que la pagina (que
% haria imposible colocar el flotante).
% Si el archivo no existe todavia, dibuja un recuadro con el nombre que
% falta en lugar de detener la compilacion.
\newcommand{\altomaximagen}{0.80\textheight}
\newcommand{\imagen}[2]{%
  \IfFileExists{#2}%
    {\includegraphics[width=#1,height=\altomaximagen,keepaspectratio]{#2}}%
    {\fbox{\parbox[c][30mm][c]{\dimexpr#1-2\fboxsep-2\fboxrule\relax}%
      {\centering\footnotesize\textsf{[Falta la imagen]\\\texttt{#2}}}}}%
}

%=======================================================================
%  3. TITULO Y AUTORES
%=======================================================================
\title{\bajartitulo Simulación de un automatismo eléctrico mediante diagrama
de escalera (ladder) en CADe SIMU}

\author{
\IEEEauthorblockN{Nombre Apellido Apellido}
\IEEEauthorblockA{Código: 000000000\\
\programa\\
\universidad\\
correo1@universidad.edu}
\and
\IEEEauthorblockN{Nombre Apellido Apellido}
\IEEEauthorblockA{Código: 000000000\\
\programa\\
\universidad\\
correo2@universidad.edu}
\and
\IEEEauthorblockN{Nombre Apellido Apellido}
\IEEEauthorblockA{Código: 000000000\\
\programa\\
\universidad\\
correo3@universidad.edu}
}

%=======================================================================
\begin{document}

% Etiquetas al estilo IEEE (van despues de \begin{document} para
% que no las sobreescriba babel)
\renewcommand{\figurename}{Fig.}
\renewcommand{\tablename}{TABLA}

\encabezadoinstitucional   % <- dibuja el logo en el margen de la pagina 1
\maketitle
\thispagestyle{plain}      % numero de pagina tambien en la primera hoja
\pagestyle{plain}          % comenta estas dos lineas si no quieres numeracion

%-----------------------------------------------------------------------
\begin{abstract}
% RESUMEN: 150-250 palabras. Escribelo AL FINAL, cuando ya sepas que dio
% la simulacion. Debe responder: que se simulo, con que herramienta, como
% se verifico y cual fue el resultado principal.
En este informe se documenta el diseño, la implementación y la validación
de un automatismo eléctrico para el control de un motor asíncrono trifásico,
desarrollado en lenguaje de contactos (diagrama ladder) y simulado en la
herramienta CADe SIMU. Se describe el proceso a automatizar, la asignación de
entradas y salidas, las ecuaciones lógicas que gobiernan el circuito de mando
y la construcción de los circuitos de fuerza y control. La verificación se
realizó mediante un protocolo de pruebas que contrasta cada requisito
funcional con el estado observado en la simulación, incluyendo las condiciones
de falla y las maniobras de protección. Los resultados muestran que el
automatismo cumple XX de XX requisitos definidos, y se discuten las
diferencias encontradas entre el comportamiento simulado y el esperado.
\end{abstract}

\begin{IEEEkeywords}
Automatismo eléctrico, diagrama de escalera, CADe SIMU, contactor,
enclavamiento, simulación, motor trifásico.
\end{IEEEkeywords}

% --- Abstract en ingles: eliminalo si el docente no lo exige -----------
\vspace{6pt}
\noindent\textit{\textbf{Abstract}---This report documents the design,
implementation and validation of an electrical control system for a
three-phase induction motor, developed in ladder logic and simulated with
CADe SIMU. The process description, I/O assignment, logic equations, power
and control circuits are presented, together with a test protocol that
contrasts every functional requirement against the simulated behaviour,
including fault and protection conditions.}

\vspace{4pt}
\noindent\textit{\textbf{Index Terms}---Electrical control, ladder diagram,
CADe SIMU, contactor, interlocking, simulation, three-phase motor.}
\vspace{6pt}

%=======================================================================
\section{Introducción}
%=======================================================================
% Que responde esta seccion:
%   1) contexto: por que existen los automatismos cableados
%   2) problema: que se quiere controlar y por que es relevante
%   3) por que simular antes de cablear
%   4) que hace este informe y como esta organizado
\IEEEPARstart{L}{os} automatismos eléctricos cableados constituyen la base del
control industrial de maquinaria, y su representación mediante diagramas de
escalera o \textit{ladder} sigue siendo el lenguaje más extendido tanto en
lógica cableada como en la programación de controladores lógicos programables,
normalizado en la IEC~61131-3 \cite{iec61131}. Un error en el circuito de mando
puede provocar cortocircuitos, daño de equipos o riesgo para el operario, por
lo que la simulación previa al montaje es una etapa necesaria del diseño.

CADe SIMU es una herramienta de diseño y simulación de circuitos eléctricos
que permite dibujar los esquemas de fuerza y de mando con simbología
normalizada y ejecutar la simulación observando en tiempo real qué elementos
conducen \cite{cadesimu}. Esto permite validar la lógica antes de llevarla al
tablero real.

En este trabajo se implementa y simula \textbf{[describe en una frase el
automatismo: p.\,ej. el arranque directo con inversión de giro de un motor
trifásico]}. El documento se organiza así: la Sección~\ref{sec:objetivos}
plantea los objetivos; la Sección~\ref{sec:marco} resume los fundamentos
teóricos; la Sección~\ref{sec:materiales} lista los materiales y el software;
la Sección~\ref{sec:proceso} describe el proceso a automatizar y sus
requisitos; la Sección~\ref{sec:desarrollo} detalla el desarrollo de la
simulación; la Sección~\ref{sec:resultados} presenta las pruebas y resultados;
la Sección~\ref{sec:analisis} los analiza y la Sección~\ref{sec:conclusiones}
concluye.

%=======================================================================
\section{Objetivos}\label{sec:objetivos}
%=======================================================================
\subsection{Objetivo general}
Diseñar, implementar y validar mediante simulación en CADe SIMU un automatismo
eléctrico basado en diagrama de escalera que cumpla las condiciones de
operación y protección definidas para \textbf{[el proceso]}.

\subsection{Objetivos específicos}
\begin{itemize}[leftmargin=1.2em,itemsep=1pt,topsep=2pt]
  \item Identificar los elementos de mando, maniobra, protección y
        señalización requeridos por el proceso.
  \item Definir la tabla de asignación de entradas y salidas del automatismo.
  \item Obtener las ecuaciones lógicas del circuito de mando y traducirlas a
        lenguaje de contactos.
  \item Construir los circuitos de fuerza y de control en CADe SIMU con
        simbología normalizada.
  \item Verificar mediante un protocolo de pruebas el cumplimiento de cada
        requisito funcional, incluidas las condiciones de falla.
\end{itemize}

%=======================================================================
\section{Marco teórico}\label{sec:marco}
%=======================================================================
% Regla practica: solo teoria que uses despues. Si un concepto no aparece
% en el desarrollo o en el analisis, sobra aqui.

\subsection{Diagrama de escalera}
El diagrama de escalera representa el circuito de mando entre dos líneas
verticales de alimentación; cada renglón o \textit{peldaño} contiene la
combinación de contactos (condiciones) que energiza una bobina o salida
\cite{petruzella}. Los contactos normalmente abiertos (NA) representan la
función lógica directa y los normalmente cerrados (NC) su complemento, de modo
que un renglón equivale a una función booleana de sus entradas.

\subsection{Elementos de mando, maniobra y protección}
% Describe brevemente SOLO los que usaste: pulsadores NA/NC, selector,
% contactor, rele termico, guardamotor, fusibles, lamparas de senalizacion,
% final de carrera, seta de emergencia...
Los pulsadores de marcha y paro, el contactor como elemento de maniobra, el
relé térmico y el guardamotor como protecciones, y las lámparas de
señalización conforman el conjunto mínimo de un automatismo de arranque. La
simbología empleada corresponde a la norma IEC~60617 \cite{iec60617}.

\subsection{Realimentación y enclavamientos}
El contacto auxiliar de la propia bobina en paralelo con el pulsador de marcha
constituye la \textit{realimentación} o enclavamiento eléctrico, que mantiene
el circuito energizado tras soltar el pulsador. El \textit{enclavamiento
mutuo} entre dos contactores se implementa intercalando el contacto NC de cada
uno en el renglón del otro, e impide maniobras incompatibles como la
alimentación simultánea de dos sentidos de giro \cite{roldan}.

\subsection{Temporizadores y contadores}
% Elimina esta subseccion si tu practica no los usa.
El temporizador a la conexión (TON) retarda la activación de su contacto un
tiempo $t$ tras energizarse su bobina, mientras que el temporizador a la
desconexión (TOFF) retarda la desactivación. Son la base de secuencias
automáticas como el arranque estrella-triángulo.

\subsection{CADe SIMU como entorno de simulación}
CADe SIMU permite dibujar sobre una hoja los circuitos de fuerza y mando
empleando librerías de símbolos normalizados y ejecutar la simulación, durante
la cual los conductores energizados cambian de color y los contactos asociados
a una bobina conmutan automáticamente \cite{cadesimu}. La herramienta no
resuelve el comportamiento eléctrico transitorio: valida la \textit{lógica} de
la maniobra y la continuidad del circuito, lo que debe tenerse presente al
interpretar los resultados.

%=======================================================================
\section{Materiales y software}\label{sec:materiales}
%=======================================================================
En la Tabla~\ref{tab:materiales} se listan los elementos empleados en la
simulación con su designación en el esquema.

\begin{table}[!t]
\renewcommand{\arraystretch}{1.15}
\caption{Elementos empleados en la simulación}
\label{tab:materiales}
\centering
\footnotesize
\begin{tabular}{@{}clcl@{}}
\toprule
\textbf{Ítem} & \textbf{Elemento} & \textbf{Cant.} & \textbf{Designación} \\
\midrule
1 & Motor asíncrono trifásico, 220/380 V & 1 & M1 \\
2 & Contactor tripolar + aux.            & 2 & KM1, KM2 \\
3 & Relé térmico                         & 1 & F1 \\
4 & Guardamotor / magnetotérmico         & 1 & Q1 \\
5 & Pulsador NA (marcha)                 & 2 & S1, S2 \\
6 & Pulsador NC (paro)                   & 1 & S0 \\
7 & Seta de emergencia con retención     & 1 & S3 \\
8 & Lámpara de señalización              & 3 & H1, H2, H3 \\
\bottomrule
\end{tabular}
\end{table}

\textbf{Software:} CADe SIMU versión \textbf{[X.X]}, ejecutado sobre
\textbf{[sistema operativo]}. El archivo de simulación se entrega como anexo
con el nombre \texttt{[nombre\_archivo].dsm}.

%=======================================================================
\section{Descripción del proceso a automatizar}\label{sec:proceso}
%=======================================================================
% Aqui va el enunciado de la practica REDACTADO POR USTEDES, no copiado.
% Separa la narracion (que hace la maquina) de los requisitos (lista
% numerada y verificable). Esos requisitos son los que vas a probar
% uno por uno en la Seccion de resultados: por eso se numeran.

\subsection{Planteamiento}
\textbf{[Describe la máquina o el proceso: qué acciona el motor, en qué
condiciones arranca y se detiene, qué debe señalizarse, qué ocurre ante una
falla.]}

\subsection{Requisitos funcionales}
\begin{enumerate}[label=\textbf{R\arabic*.},leftmargin=2.2em,itemsep=1pt,topsep=2pt]
  \item Al accionar S1 el motor arranca en sentido directo y permanece en
        marcha al soltar el pulsador.
  \item Al accionar S2 el motor arranca en sentido inverso bajo las mismas
        condiciones.
  \item No debe ser posible energizar KM1 y KM2 simultáneamente.
  \item El pulsador S0 detiene el motor en cualquier estado.
  \item El disparo del relé térmico F1 detiene el motor y activa la
        señalización de falla H3.
  \item La lámpara H1 indica marcha y H2 indica paro.
\end{enumerate}

%=======================================================================
\section{Desarrollo de la simulación}\label{sec:desarrollo}
%=======================================================================
% Esta es la seccion que mas peso tiene. El orden propuesto va de lo
% abstracto a lo concreto: variables -> logica -> esquema -> procedimiento.

\subsection{Asignación de entradas y salidas}
La Tabla~\ref{tab:io} relaciona cada elemento físico con la variable lógica
empleada en las ecuaciones y en el diagrama de escalera.

\begin{table}[!t]
\renewcommand{\arraystretch}{1.15}
\caption{Asignación de entradas y salidas}
\label{tab:io}
\centering
\footnotesize
\begin{tabular}{@{}llll@{}}
\toprule
\textbf{Variable} & \textbf{Elemento} & \textbf{Tipo} & \textbf{Función} \\
\midrule
S0 & Pulsador NC   & Entrada & Paro general \\
S1 & Pulsador NA   & Entrada & Marcha directa \\
S2 & Pulsador NA   & Entrada & Marcha inversa \\
F1 & Contacto NC del relé térmico & Entrada & Protección \\
KM1 & Bobina de contactor & Salida & Giro directo \\
KM2 & Bobina de contactor & Salida & Giro inverso \\
H1 & Lámpara & Salida & Señalización de marcha \\
H3 & Lámpara & Salida & Señalización de falla \\
\bottomrule
\end{tabular}
\end{table}

\subsection{Ecuaciones lógicas del circuito de mando}
A partir de los requisitos se obtienen las ecuaciones booleanas de cada
bobina. La realimentación aparece como el término $KM_i$ en paralelo con el
pulsador de marcha, y el enclavamiento mutuo como los factores
$\overline{KM_2}$ y $\overline{KM_1}$:
\begin{align}
  KM_1 &= (S_1 + KM_1)\cdot \overline{S_0}\cdot \overline{F_1}\cdot \overline{KM_2} \label{eq:km1}\\
  KM_2 &= (S_2 + KM_2)\cdot \overline{S_0}\cdot \overline{F_1}\cdot \overline{KM_1} \label{eq:km2}\\
  H_1  &= KM_1 + KM_2 \label{eq:h1}\\
  H_3  &= F_1 \label{eq:h3}
\end{align}
Cada ecuación se traduce directamente a un renglón del diagrama de escalera:
los términos en producto se conectan en serie y los términos en suma en
paralelo.

\subsection{Circuito de fuerza}
En la Fig.~\ref{fig:fuerza} se muestra el circuito de fuerza, con la
protección Q1, los contactos principales de los contactores, el relé térmico
F1 y el motor M1. La inversión de giro se obtiene permutando dos de las tres
fases en KM2.

\begin{figure}[!t]
  \centering
  % Sustituye por la captura de CADe SIMU del circuito de fuerza
  \imagen{\columnwidth}{figuras/circuito-fuerza.png}
  \caption{Circuito de fuerza implementado en CADe SIMU.}
  \label{fig:fuerza}
\end{figure}

\subsection{Circuito de control en lenguaje de contactos}
La Fig.~\ref{fig:ladder} presenta el diagrama de escalera resultante de las
ecuaciones (\ref{eq:km1})--(\ref{eq:h3}). Se emplea una figura a dos columnas
para que los renglones sean legibles; si tu esquema es angosto, usa el entorno
\texttt{figure} normal como en la Fig.~\ref{fig:fuerza}.

\begin{figure*}[!t]
  \centering
  % Esta figura ocupa TODO EL ANCHO DEL PAPEL: la caja mide \textwidth y la
  % imagen es mas ancha, de modo que sobresale por igual a los dos margenes.
  % Variantes:
  %   solo el ancho de las dos columnas -> \imagen{\textwidth}{...} sin \makebox
  %   una sola columna                  -> figure (sin *) e \imagen{\columnwidth}{...}
  \hfuzz=30mm  % la imagen se sale a proposito: silencia el aviso Overfull
  \makebox[\textwidth][c]{\imagen{0.95\paperwidth}{figuras/diagrama-ladder.png}}
  \caption{Diagrama de escalera del circuito de mando. Cada renglón
  corresponde a una de las ecuaciones lógicas planteadas.}
  \label{fig:ladder}
\end{figure*}

\subsection{Procedimiento seguido en CADe SIMU}
% Escribe lo que HICIERON, en pasado y en pasos reproducibles. Un tercero
% deberia poder repetir la simulacion leyendo esto.
\begin{enumerate}[leftmargin=1.6em,itemsep=1pt,topsep=2pt]
  \item Se creó un nuevo esquema y se seleccionó la norma de simbología
        \textbf{[IEC/otra]} y el número de hojas requerido.
  \item Se trazó la alimentación trifásica y se insertaron Q1, los contactos
        principales de KM1 y KM2, F1 y el motor M1.
  \item Se dibujó el circuito de mando entre las líneas de alimentación,
        insertando pulsadores, contactos auxiliares, bobinas y lámparas.
  \item Se etiquetó cada elemento con la designación de la
        Tabla~\ref{tab:io}, de modo que los contactos auxiliares quedaran
        asociados a su bobina.
  \item Se ejecutó la simulación y se accionaron los elementos de mando
        siguiendo el protocolo de pruebas de la Tabla~\ref{tab:pruebas},
        registrando el estado de las salidas en cada caso.
\end{enumerate}

%=======================================================================
\section{Pruebas y resultados}\label{sec:resultados}
%=======================================================================
% Resultados = lo que se OBSERVO, sin interpretarlo todavia.
% Interpretar es tarea de la seccion siguiente.

\subsection{Protocolo de pruebas}
Cada requisito de la Sección~\ref{sec:proceso} se verificó accionando la
entrada correspondiente y observando el estado de las salidas durante la
simulación. La Tabla~\ref{tab:pruebas} resume los casos ejecutados.

\begin{table*}[!t]
\renewcommand{\arraystretch}{1.2}
\caption{Protocolo de pruebas y resultados observados en la simulación}
\label{tab:pruebas}
\centering
\footnotesize
\begin{tabular}{@{}clllcc@{}}
\toprule
\textbf{\#} & \textbf{Req.} & \textbf{Acción ejecutada} & \textbf{Resultado esperado} &
\textbf{Resultado observado} & \textbf{Cumple} \\
\midrule
1 & R1 & Pulsar y soltar S1 & KM1 energizado, se mantiene; H1 encendida & & Sí / No \\
2 & R2 & Pulsar y soltar S2 con motor parado & KM2 energizado, se mantiene & & Sí / No \\
3 & R3 & Pulsar S2 con KM1 activo & KM2 no se energiza & & Sí / No \\
4 & R4 & Pulsar S0 en marcha & Ambas bobinas se desenergizan & & Sí / No \\
5 & R5 & Forzar disparo de F1 & Motor se detiene; H3 encendida & & Sí / No \\
6 & R6 & Observar señalización en reposo & H2 encendida, H1 apagada & & Sí / No \\
\bottomrule
\end{tabular}
\end{table*}

\subsection{Estados de la simulación}
Las Figs.~\ref{fig:estado-reposo} y \ref{fig:estado-marcha} muestran el estado
del circuito en reposo y en marcha directa; los conductores energizados se
distinguen por el cambio de color que aplica el simulador.

\begin{figure}[!t]
  \centering
  \imagen{\columnwidth}{figuras/estado-reposo.png}
  \caption{Estado del automatismo en reposo: ninguna bobina energizada.}
  \label{fig:estado-reposo}
\end{figure}

\begin{figure}[!t]
  \centering
  \imagen{\columnwidth}{figuras/estado-marcha.png}
  \caption{Estado del automatismo tras accionar S1: KM1 energizado y
  realimentado por su contacto auxiliar.}
  \label{fig:estado-marcha}
\end{figure}

\subsection{Diagrama de tiempos}
La Fig.~\ref{fig:cronograma} representa la secuencia temporal de las señales
durante una maniobra completa de marcha y paro.
% Ajusta los quiebres de cada linea a TU secuencia. Coordenada x = tiempo.
\begin{figure}[!t]
\centering
\begin{tikzpicture}[x=0.85cm,y=0.5cm,line width=0.7pt]
  \scriptsize
  % Instantes de conmutacion
  \draw[dashed,gray!70] (1,-0.6) -- (1,7.2);
  \draw[dashed,gray!70] (1.8,-0.6) -- (1.8,7.2);
  \draw[dashed,gray!70] (5,-0.6) -- (5,7.2);
  \node[gray!80] at (1,7.6) {$t_1$};
  \node[gray!80] at (1.8,7.6) {$t_2$};
  \node[gray!80] at (5,7.6) {$t_3$};
  % S1 (marcha): pulso t1..t2
  \node[anchor=east] at (-0.1,6) {S1};
  \draw (0,6) -- (1,6) -- (1,7) -- (1.8,7) -- (1.8,6) -- (7,6);
  % S0 (paro, NC): pulso en t3
  \node[anchor=east] at (-0.1,4) {S0};
  \draw (0,4) -- (5,4) -- (5,5) -- (5.6,5) -- (5.6,4) -- (7,4);
  % KM1: se energiza en t1 y cae en t3
  \node[anchor=east] at (-0.1,2) {KM1};
  \draw (0,2) -- (1,2) -- (1,3) -- (5,3) -- (5,2) -- (7,2);
  % H1: sigue a KM1
  \node[anchor=east] at (-0.1,0) {H1};
  \draw (0,0) -- (1,0) -- (1,1) -- (5,1) -- (5,0) -- (7,0);
  % Eje de tiempo
  \draw[->] (0,-0.6) -- (7.3,-0.6) node[right] {$t$};
\end{tikzpicture}
\caption{Diagrama de tiempos de la maniobra de marcha directa y paro.}
\label{fig:cronograma}
\end{figure}

%=======================================================================
\section{Análisis de resultados}\label{sec:analisis}
%=======================================================================
% Aqui SI se interpreta. Preguntas que conviene responder:
%   - Coincide lo observado con lo esperado? Donde no, y por que?
%   - Que errores aparecieron al construir el esquema y como se corrigieron?
%   - Que papel jugo cada enclavamiento: que pasaria si se quitara?
%   - Que limites tiene la simulacion frente al montaje real?
\textbf{[Contrasta cada fila de la Tabla~\ref{tab:pruebas} con lo esperado.]}
El enclavamiento mutuo introducido por los factores $\overline{KM_2}$ y
$\overline{KM_1}$ en (\ref{eq:km1}) y (\ref{eq:km2}) resultó determinante: al
retirarlo de forma experimental, la simulación permitió energizar ambas
bobinas simultáneamente, condición que en el circuito de fuerza equivale a un
cortocircuito entre fases.

Debe señalarse que CADe SIMU valida la continuidad y la lógica de la maniobra,
pero no modela el tiempo de conmutación de los contactos ni el arco eléctrico,
de modo que fenómenos como el rebote de contactos o la necesidad de un
enclavamiento mecánico adicional entre contactores no se evidencian en la
simulación y deben considerarse en el montaje real.

%=======================================================================
\section{Conclusiones}\label{sec:conclusiones}
%=======================================================================
% Una conclusion por objetivo especifico. Deben responder a lo que
% SE HIZO, no a lo que se aprendio en general.
\begin{itemize}[leftmargin=1.2em,itemsep=1pt,topsep=2pt]
  \item \textbf{[Conclusión sobre el cumplimiento de los requisitos.]}
  \item \textbf{[Conclusión sobre la traducción ecuaciones $\rightarrow$ ladder.]}
  \item \textbf{[Conclusión sobre el papel de las protecciones y enclavamientos.]}
  \item \textbf{[Conclusión sobre la utilidad y los límites de la simulación.]}
\end{itemize}

%=======================================================================
\section{Recomendaciones}
%=======================================================================
% Opcional: borra la seccion si no la piden.
\textbf{[Mejoras al automatismo, precauciones para el montaje real,
trabajo futuro.]}

%=======================================================================
\appendices
\section{Aportes de los integrantes}
\begin{table}[!h]
\renewcommand{\arraystretch}{1.15}
\centering
\footnotesize
\begin{tabular}{@{}ll@{}}
\toprule
\textbf{Integrante} & \textbf{Aporte principal} \\
\midrule
Nombre 1 & Diseño de la lógica y ecuaciones \\
Nombre 2 & Montaje del esquema en CADe SIMU \\
Nombre 3 & Protocolo de pruebas y redacción \\
\bottomrule
\end{tabular}
\end{table}

\section{Archivos anexos}
Se adjunta el archivo de simulación \texttt{[nombre\_archivo].dsm} y las
capturas en resolución original empleadas en las figuras de este informe.

%=======================================================================
\bibliographystyle{IEEEtran}
\bibliography{referencias}

\end{document}
